\bookchapter{第五十三章：資料儲存與網路請求}

\chapterintro{一個只能“看”的 App 是死的；能存資料、能連網路的 App 才是活的。這一章認識安卓資料層的“三件套”：DataStore（存小配置）、Room（存結構化資料）、Retrofit（和伺服器對話），外加它們的“總排程”——協程。學完這章，你就能寫會記住使用者習慣、會實時更新內容的真 App 了。}

\booksection{一、儲存方案怎麼選：先看這張表}

\begin{center}
\begin{tabularx}{\linewidth}{|X|X|X|X|}
\hline
\textbf{方案} & \textbf{適合存什麼} & \textbf{特點} & \textbf{本書哪裡用} \tabularnewline
\hline
DataStore & 零散配置（主題、開關） & 鍵值對，簡單安全 & 第 54 章記帳本設定 \tabularnewline
Room & 結構化資料（賬目、使用者） & 真正的 SQLite 資料庫封裝，可查詢 & 第 54 章賬目記錄 \tabularnewline
Retrofit & （不存，聯網取）伺服器資料 & HTTP 請求庫的事實標準 & 第 55 章天氣資料 \tabularnewline
\hline
\end{tabularx}
\end{center}

\begin{center}\small\color{gray}資料三件套的分工\end{center}

記一個口訣：“小配置進 DataStore，成批資料進 Room，外面的資料走 Retrofit”。舊教學裡常見的 SharedPreferences 已被 DataStore 取代，SQLite 手寫 SQL 也已被 Room 取代——直接學新的，別走老路。

\booksection{二、DataStore：存點小配置}

想記住“使用者上次選的幣種”“是否開啟深色模式”這類零散值，用 DataStore（Preferences 版）：

\begin{lstlisting}[language=Kotlin]
val Context.dataStore by preferencesDataStore(name = "settings")

// 寫：存一個值
suspend fun saveTheme(isDark: Boolean) {
    context.dataStore.edit { prefs ->
        prefs[THEME_KEY] = isDark
    }
}

// 讀：取出來（Flow 會自動通知變化）
val themeFlow: Flow<Boolean> = context.dataStore.data
    .map { it[THEME_KEY] ?: false }
\end{lstlisting}

注意返回的是 Flow——一個會“流動”的資料流。值一旦改變，所有觀察它的地方自動收到新值。這個概念（響應式）是安卓現代的基石，第 54 章會實戰。

\booksection{三、Room：給資料安個“真正的家”}

記帳本有成百上千筆賬，還要按分類、按月份查詢統計——這需要資料庫。Room 是 SQLite 的官方封裝，把“建表、查資料”變得像寫 Kotlin 一樣自然。三樣東西：

\begin{lstlisting}[language=Kotlin]
// 1. Entity：表結構（一張表 = 一個 data class）
@Entity(tableName = "records")
data class Record(
    @PrimaryKey(autoGenerate = true) val id: Long = 0,
    val amount: Double,
    val category: String,
    val note: String,
    val date: String
)

// 2. Dao：操作介面（增刪改查）
@Dao
interface RecordDao {
    @Query("SELECT * FROM records ORDER BY date DESC")
    fun getAll(): Flow<List<Record>>

    @Insert
    suspend fun insert(record: Record)

    @Delete
    suspend fun delete(record: Record)

    @Query("SELECT SUM(amount) FROM records WHERE category = :c")
    suspend fun sumByCategory(c: String): Double
}

// 3. Database：總裝
@Database(entities = [Record::class], version = 1)
abstract class AppDatabase : RoomDatabase() {
    abstract fun recordDao(): RecordDao
}
\end{lstlisting}

這套“三件套”模板可以原樣搬進你的任何專案——換換欄位名就行。查詢返回 Flow，介面跟著資料自動重新整理，妙不可言。

\booksection{四、Retrofit：和伺服器說話}

天氣、新聞、匯率……這些資料不在手機裡，在別人的伺服器上。Retrofit 讓“向伺服器要資料”變得像調函式一樣簡單。三步走：

\begin{lstlisting}[language=Kotlin]
// 1. 定義介面：把 HTTP 請求描述成函式
interface WeatherApi {
    @GET("v1/weather")
    suspend fun getWeather(@Query("city") city: String): WeatherResponse
}

// 2. 建立例項
val api = Retrofit.Builder()
    .baseUrl("https://api.example.com/")
    .addConverterFactory(GsonConverterFactory.create())
    .build()
    .create(WeatherApi::class.java)

// 3. 像調函式一樣呼叫（在協程裡）
val result = api.getWeather("北京")
\end{lstlisting}

\textbf{避坑提醒：兩條網路鐵律}\par

(1) 網路請求絕不能在主執行緒做——會卡死介面甚至崩潰（NetworkOnMainThreadException）。所有網路呼叫必須放在協程（viewModelScope.launch）裡。(2) 明文 HTTP 預設被禁止：安卓 9 起要求 HTTPS。測試時若要用 http:// 的本地地址，需要在 Manifest 裡顯式宣告 usesCleartextTraffic——但正式產品必須全走 HTTPS。

\booksection{五、JSON 解析：給網路資料“翻譯”}

伺服器返回的資料大多是 JSON——一串像“巢狀字典”的文字。Retrofit 需要一個“翻譯官”把它變成 Kotlin 物件。2026 年的兩大主流方案：

\begin{center}
\begin{tabularx}{\linewidth}{|X|X|X|}
\hline
\textbf{方案} & \textbf{特點} & \textbf{推薦場景} \tabularnewline
\hline
Gson / Moshi & 老牌反射方案，配置零成本 & 快速起步、欄位不多的小專案 \tabularnewline
kotlinx.serialization & 官方編譯期方案，更快更安全 & 新專案首選（2026 年主流） \tabularnewline
\hline
\end{tabularx}
\end{center}

\begin{center}\small\color{gray}JSON 解析兩大方案\end{center}

\begin{lstlisting}[language=Kotlin]
// 給資料類加 @Serializable 註解
@Serializable
data class WeatherResponse(
    val weather: List<Weather>,
    val main: Main,
    val name: String
)

// Retrofit 掛上對應 Converter（以 Retrofit 2.11+ 官方支援為例）：
// .addConverterFactory(json.asConverterFactory("application/json".toMediaType()))
\end{lstlisting}

無論哪種方案，鐵律都是：先看真實返回的 JSON，再定義欄位。很多“解析失敗”的 bug，都是因為欄位名寫錯、巢狀層級猜錯——把原始 JSON 列印出來對照，十秒定位。

\booksection{六、協程：安卓併發的“普通話”}

\booksection{八、協程：安卓併發的“普通話”}

上面反覆出現的 suspend 和 launch 來自 Kotlin 協程（Coroutine）——安卓官方的併發方案。你可以粗暴地理解為：讓“要等一會兒”的操作（網路、資料庫）不阻塞介面。

\begin{lstlisting}[language=Kotlin]
// 在 ViewModel 裡啟動協程
viewModelScope.launch {
    val records = dao.getAll().first()  // suspend：會自動掛起等待
    // 更新介面狀態……
}

// 關鍵字速記：
// suspend   —— "我要等待，別攔著別人"
// launch    —— "開工！"
// viewModelScope —— 跟著頁面走的"工作間"，頁面銷燬自動收工
\end{lstlisting}

零基礎先記住用法就好，背後的“掛起與還原”原理可以以後再深究。日常開發你寫的 90\% 都是這個模板。

\begin{tcbpractice}{來動手：做一個"每日一句"小應用}
找一個免費的名言 API（比如 hitokoto 之類），用 Retrofit 拉取“每日一句”，顯示在螢幕中央，加個“換一句”按鈕（協程裡重新請求）。再把喜歡的句子用 Room 存進“收藏夾”。這個練習把 DataStore/Room/Retrofit/協程全用了一遍——做完它，第 7、8 章就是小菜一碟。

\end{tcbpractice}

\begin{tcbsummary}{本章小結}
\begin{itemize}
  \item 口訣：小配置進 DataStore，成批資料進 Room，外面的資料走 Retrofit
  \item DataStore 返回 Flow：值變了，觀察者自動收到
  \item Room 三件套：@Entity（表）+ @Dao（操作）+ @Database（總裝），模板可直接複用
  \item Retrofit 三步：定義介面 建立例項 協程裡呼叫
  \item 兩條鐵律：網路不進主執行緒；正式產品必須 HTTPS
  \item 協程 = 不阻塞介面的等待；viewModelScope 是頁面專屬工作間
\end{itemize}

\end{tcbsummary}
